\section{Moment losses and a common scale}\label{sec:shape}
This section proves the approximate core used in both main theorems.
We normalize $B=1$ and omit row indices. Thus $R=L$. Put
$v_i=\E X_i^2$, $\sigma_i=\sqrt{v_i}$, $\rho_i=\E|X_i|^3$,
$\kappa_i=\E X_i^3$, and $b_i=\E|X_i|$. Define
\[
 \delta_i=\sqrt{10}\rho_i-|\kappa_i|-3b_iv_i.
\]

\begin{lemma}[Retained losses]\label{lem:retained}
Under \eqref{eq:near-assumption},
\begin{equation}\label{eq:retained-losses}
 \sum_i\delta_i=o(R),\quad
 \sum_i\E[|X_i|^3\ind_{\{|X_i|>R^{1/6}\}}]=o(R),\quad
 \sum_i|\kappa_i|-\Big|\sum_i\kappa_i\Big|=o(R),
\end{equation}
and $\sum_i\sigma_i^3/R\to1/\be$. After one reflection of a row,
let $P_i^*$ be the Esseen two-point law of variance $v_i$ with the
same positive orientation for every $i$. Then
\begin{equation}\label{eq:shape-W3}
                  \sum_iW_3^3(X_i,P_i^*)=o(R).
\end{equation}
\end{lemma}
\begin{proof}
First we record the equality and stability part of the scalar moment
inequality. Put $s=\sqrt{\pe\qe}$, $a=\sqrt{\qe/\pe}$,
$b=\sqrt{\pe/\qe}$ and $t=\te$. For a centered variance-one law
oriented to have nonnegative third moment, set
\[
 Q(x)=-6sx^2+6tx+6s,\qquad
 f(x)=x^3-\sqrt{10}|x|^3+3|x|.
\]
Direct factorization gives
\[
 Q(x)-f(x)=
 \begin{cases}
 (\sqrt{10}-1)(x-a)^2(x+2s),&x\ge0,\\
 (\sqrt{10}+1)(-x-b)^2(-x+2s),&x\le0.
 \end{cases}
\]
Since $\E Q(X)=0$, $\delta(X)=\E(Q-f)\ge0$; equality forces
the two Esseen support points, with their centered probabilities.
This also checks the equality statement needed from
\citet[Theorem~3.1, pp.~256--262]{Shevtsova2012}.

For $c_0=2\pe(\sqrt{10}-1)>0$, each factor is at least $c_0$
times the cube of the distance to its half-line contact. Round a
variance-one law to $a$ on the positive half-line and to $-b$ on the
negative half-line. If $e$ is its cubic rounding cost, then
$e\le\delta(X)/c_0$ and centering gives
$|\Pp\{\hbox{rounded value}=a\}-\pe|\le e^{1/3}/(a+b)$.
Couple the rounded law to the exact Esseen probabilities. The inequality
$(u+v)^3\le4(u^3+v^3)$ yields
\[
 W_3^3(X,W_+)\le4e+4(a+b)^2e^{1/3}.
\]
After rescaling coordinates and applying H\"older,
\begin{equation}\label{eq:shape-stability}
 \frac{\sum_iW_3^3(X_i,P_i^{\rm oriented})}{R}
 \le\frac4{c_0}\bar\delta+
       4(a+b)^2c_0^{-1/3}\bar\delta^{1/3},\qquad
 \bar\delta=R^{-1}\sum_i\delta_i.
\end{equation}
Here the orientations initially follow the signs of individual
$\kappa_i$; zero laws have zero cost.

The proof on pp.~289--296 of \citet{Shevtsova2012} permits retention
of its nonnegative losses. We give the details of this extraction.
Put $U=R^{1/6}$, $Y_i=X_i\ind_{\{|X_i|\le U\}}$,
$Z_i=Y_i-\E Y_i$, and $T_i=\E[|X_i|^3\ind_{\{|X_i|>U\}}]$.
The source's $I_{14},I_{15},I_{16}$ comparisons, with its truncation
parameter $u=1$, and Lemma~4.3 on pp.~272--274 give
\[
 \Delta\le\phz(I_{16}+R/2)+O(R^{7/6}),
\]
\[
 I_{16}=\frac{\sqrt{10}R}{6}-\frac16\sum_i\delta(Z_i)
                   +I_{17}-I_{18},
\]
\[
 I_{17}=\frac{\sqrt{10}}6\sum_i(\E|Z_i|^3-\E|Y_i|^3),\qquad
 I_{18}=\frac{\sqrt{10}-1}{6}\sum_iT_i.
\]
The second identity is the moment decomposition before discarding
its deficit. At the parameter $\sqrt{10}$ the source's moment
remainder is zero. All remaining terms are of order $R^{7/6}$ or
smaller, uniformly over the independent row.

Let $b_{0i}=\rho_i^{1/3}/U\le R^{1/6}$. The bound
$|\E Y_i|\le T_i/U^2$ and cubic Taylor expansion show
\[
 \E|Z_i|^3-\E|Y_i|^3
       \le(3b_{0i}^2+3b_{0i}^4+b_{0i}^6)T_i.
\]
For $R\le1/4096$, $b_{0i}\le1/4$ and the bracket is less than
$1/3$. Since $\sqrt{10}/3<(\sqrt{10}-1)/2$, the displayed
decomposition implies
\[
 \Delta\le\ce R-\frac{\phz}{6}\sum_i\delta(Z_i)
            -\frac{\phz(\sqrt{10}-1)}{12}\sum_iT_i+O(R^{7/6}).
\]
Near-Esseen competition forces both losses to be $o(R)$. Cubic Taylor
bounds give
$|\rho_i-\rho(Z_i)|+|\kappa_i-\kappa(Z_i)|\le K T_i$, while
$\E|Z_i|\Var Z_i-\E|X_i|\Var X_i\le b_{0i}^2T_i$.
Consequently $\delta_i\le\delta(Z_i)+K'T_i$, proving the first two
claims in \eqref{eq:retained-losses} and the individually oriented
version of \eqref{eq:shape-W3}.

To retain orientation we use the aggregate low-frequency integral,
before the source's individual cubic triangle inequality. Write
$K_3=\sum_i\kappa_i$, $\nu=1+\sum_i b_iv_i/R\in[1,2]$ and
$T=2\pi/(\nu R)$. On each fixed compact frequency interval the
just-proved tail estimate gives
\[
 \chi(u)=e^{-u^2/2}(1-iK_3u^3/6)+o(R),
\]
uniformly there. Indeed the summed Taylor remainders are at most
$K[UR|u|^4+\sum_iT_i|u|^3]=o(R)$ and the summed logarithmic
quadratic errors are $O(\sum_i v_i^2+\sum_i\rho_i^2)=o(R)$.

We use Lemma~4.1 and (4.4)--(4.7) on pp.~263--267 of the source.
Its smoothing kernel is
\[
 K(t)=\frac{1-|t|}{2}+\frac{i}{2}
       \left\{(1-|t|)\cot(\pi t)+\frac{\operatorname{sign}t}{\pi}\right\},
                                      \qquad0<|t|<1.
\]
The constants $\kappa_1$ and $t_1\in(1/2,1)$ below are the universal
constants in the source's (4.4)--(4.6); in particular
$t^2(1-4\pi\kappa_1t)$ is strictly positive on $0<t\le t_1$.
Choose a fixed small $t_0>0$ so that its deleted-product bound (4.7)
has Gaussian exponent at least $1/4$ for $0\le u\le t_0T$.
The one-coordinate second-order Taylor bound against a same-variance
normal gives
\[
 R^{-1}|\chi(u)-e^{-u^2/2}|\le K|u|^3e^{-u^2/4}
                           \quad(0\le u\le t_0T).
\]
The Prawitz kernel satisfies
$2|K(u/T)|/T\le1.0253/(\pi u)$ and converges to $1/(\pi u)$.
Dominated convergence therefore evaluates the source's low integral as
$\phz|K_3|/6+o(R)$. Its middle interval is exponentially small:
the function $t^2(1-4\pi\kappa_1t)$ is positive on the chosen
$[t_0,t_1]$. At the final endpoint $|K(1-u)|\le Ku$ and
$1-\cos(2\pi u)\ge8u^2$, giving $O(T^{-2})$. The Gaussian endpoint
terms are at most $\sqrt{\pi/2}/T+O(T^{-2})$. Thus
\begin{equation}\label{eq:signed-loss}
 \Delta\le\phz\left\{\frac R2+\frac12\sum_i b_iv_i
                                  +\frac{|K_3|}{6}\right\}+o(R)
 =\ce R-\frac{\phz}{6}\left\{\sum_i\delta_i+
                  \sum_i|\kappa_i|-|K_3|\right\}+o(R).
\end{equation}
This proves the orientation loss. Reflect so $K_3\ge0$. Moment
continuity under the cubic couplings gives
\[
 \sum_i|\kappa_i-\kappa(P_i^{\rm oriented})|
 \le3\left(R+\be\sum_i\sigma_i^3\right)^{2/3}
                 \left(\sum_iW_3^3(X_i,P_i^{\rm oriented})\right)^{1/3}
 =o(R).
\]
The negatively oriented target coordinates then have
$\sum\sigma_i^3=o(R)$; reflecting just those targets costs $o(R)$
in cubic transport. This proves the common-orientation assertion.

Finally \citet[Theorem~4.14, pp.~300--303]{Shevtsova2012} gives, for
each fixed $c>2\phz/3$,
\[
 \Delta\le cR+K(c)\sum_i\sigma_i^3+O_c(R^{4/3}),
\]
where its Bernoulli envelope is
\[
 K(c)=\sup_{0<p\le1/2}
 \frac{\phz(3+1-2p)/6-c\{p^2+(1-p)^2\}}{\sqrt{p(1-p)}}.
\]
At $c=\ce$ its unique maximizing parameter is $\pe$,
$K(\ce)=0$, and $K'(\ce)=-\be$. The numerator's exact maximum
and uniqueness follow either by differentiation or from the sharp
moment equality above; parameters near zero are separated since
$c>2\phz/3$. Any subsequential limit $z\in[0,1]$ of
$\sum\sigma_i^3/R$ satisfies $\ce\le c+K(c)z$ on both sides of
$\ce$. Differentiation there forces $1-\be z=0$.
\end{proof}

\begin{lemma}[An actual Fourier window]\label{lem:window}
For every fixed $\epsilon>0$, a sufficiently small fixed $t_0>0$ gives
\[
 T'=(1+\epsilon)\frac{2\pi}{\nu R},\qquad
 H=2\int_{t_0}^1|K(u)|\,|\chi(T'u)|\,du,
\]
with $\nu\to\nus$ and
\begin{equation}\label{eq:window}
 \liminf H/R\ge\frac{\phz\nus\epsilon}{2(1+\epsilon)}>0.
\end{equation}
For one fixed such window there are $\eta>0$ and $e_n\to0$ such that
simultaneously for every subset of coordinates $A$,
\begin{equation}\label{eq:subset-inverse}
             \sum_{i\in A}v_i\le\eta^{-1}
                          \frac{\sum_{i\in A}\rho_i}{R}+e_n.
\end{equation}
Moreover $\sum_iW_2^2(X_i,P_i^*)\to0$.
\end{lemma}
\begin{proof}
Common-orientation moment continuity from Lemma~\ref{lem:retained}
gives $\nu\to1+3/\sqrt{10}=\nus$ and $K_3/R\to1/\sqrt{10}$.
Repeat the source's Prawitz inequality at the enlarged cutoff $T'$,
choosing $t_0$ smaller so its low-frequency domination remains valid.
Keep the actual high integral $H$ rather than extending a previous
high-frequency upper estimate beyond its domain. The low integral is
$\phz|K_3|/6+o(R)$ and the Gaussian endpoint terms are at most
$\phz\nu R/[2(1+\epsilon)]+O(R^2)$. Near-Esseen competition gives
\eqref{eq:window}.

For any centered law and an independent copy,
\[
 |\chi_i(u)|^2\le1-v_iu^2+\tfrac43\rho_i|u|^3.
\]
This follows from $\cos x\le1-x^2/2+|x|^3/6$ and
$\E|X_i-X_i'|^3\le8\rho_i$. Independence gives for a subset
\[
 |\chi_A(u)|\le\exp(-V_Au^2/2+2R_A|u|^3/3).
\]
The fixed window lies in $[a/R,C/R]$ with $a>0$. Choose
$0<\eta\le3/(8C)$. If $R_A/V_A\le\eta R$, then the full product
on this window is at most $\exp(-a^2V_A/(4R^2))$.
The common lower bound $H\ge hR$ implies
$V_A\le e_n=(4R^2/a^2)\log(2\sup|K|/(hR))$.
Otherwise $V_A<\eta^{-1}R_A/R$. These two cases prove
\eqref{eq:subset-inverse} with one eventual index for all subsets.

Let $E_i=W_3^3(X_i,P_i^*)$, so $\sum E_i=o(R)$. Fix $0<\zeta<1$.
On the set $E_i>\zeta\rho_i$, the original price and, by
\eqref{eq:subset-inverse}, variance tend to zero. On its complement,
Minkowski gives
$\rho_i/\sigma_i^3\le\be/(1-\zeta^{1/3})^3$ and
\[
 W_2^2(X_i,P_i^*)\le E_i^{2/3}
 \le v_i\left\{\frac{\zeta\be}{(1-\zeta^{1/3})^3}\right\}^{2/3}.
\]
On the exceptional set a centered independent same-variance coupling
costs $2v_i$. Let the row tend to infinity, then let $\zeta\downarrow0$.
\end{proof}

\begin{proposition}[A variance-and-price-full Esseen core]\label{prop:core}
The IID target conclusion of Theorem~\ref{thm:near} holds. In gap units
$h=\nus R$, there is a further retained set of cardinality $M$ with
\begin{equation}\label{eq:raw-core}
 M R^2\to\be^2,\quad
 \sup_{i\in I}W_3(\Law(X_i/h),\Law(B_{\pe}-\pe))\to0,
\end{equation}
retaining $1-o(1)$ of original variance and price. In these units the
total original variance and price satisfy
$V\sim\pe\qe M$ and $P/V\to\re$, and its signed cubic sum divided
by $V$ tends to $\te$.
\end{proposition}
\begin{proof}
Consider the variance probability measures
$\mu_n=\sum_i v_i\delta_{\sigma_i/R}$.
Their first moments tend to $s_0=1/\be$, so they are tight at
infinity. They have no limiting mass at zero. To prove this, discard
the set $E_i>\zeta\rho_i$ as in the last lemma. On the remaining set
$\rho_i/\sigma_i^3\le K_\zeta$. For $\sigma_i\le\delta R$, its
price is at most $K_\zeta\delta R$ times its variance. The subset
inverse gives limiting variance at most $\eta^{-1}K_\zeta\delta$.
Portmanteau on the relative-open interval $[0,\delta)$ followed by
$\delta\downarrow0$ excludes a zero atom from every limit $\mu$.
Fatou gives $\int z\,d\mu\le s_0$; first-moment uniform
integrability has not yet been assumed.

At each fixed $\epsilon>0$, \eqref{eq:window} selects a frequency
$u_n\in[t_0T',T']$ with $|\chi(u_n)|\ge c_\epsilon R$.
Let $w_n=Ru_n$. It lies in a compact positive interval with
$\limsup w_n\le2\pi(1+\epsilon)/\nus$. Independence gives
$\sum_i(1-|\chi_i(u_n)|^2)=O_\epsilon(\log(1/R))$.
The same-variance binary target has defect
$4\pe\qe\sin^2((a+b)\sigma_i u_n/2)$.
On $\sigma_i\le KR$, the original variance-weighted defects tend
to zero, and their transfer to target defects is at most
\[
 2|u_n|\Big(\sum_{\sigma_i\le KR}v_i^2\Big)^{1/2}
                          \Big(\sum_iW_2^2(X_i,P_i^*)\Big)^{1/2}=o(1).
\]
The untruncated target contribution is bounded by $\mu_n([K,\infty))$.
First let the row tend to infinity, then $K\to\infty$. Along a weak
subsequence with $w_n\to w_\epsilon$, this proves
\[
          \int\sin^2((a+b)zw_\epsilon/2)\,d\mu(z)=0.
\]
The limit is supported on positive integer multiples of
$2\pi/((a+b)w_\epsilon)$, hence on $[s_0/(1+\epsilon),\infty)$;
we used $\be/(a+b)=1/\nus$. Repeat for $\epsilon=1/j$, taking
further subsequences of the same weak limit. That limit is supported
on $[s_0,\infty)$ and has mean at most $s_0$, so it is $\delta_{s_0}$.
The exact first moments converge; therefore
\[
 \int|z-s_0|\,d\mu_n
   =\int z\,d\mu_n-s_0+2\int(s_0-z)_+\,d\mu_n\longrightarrow0.
\]

Trim a slowly shrinking band about $s_0$ and denote its count by $m$.
It retains total variance $1-o(1)$ and target cubic resource
$\be\sum_{I}\sigma_i^3=R(1-o(1))$. The original price has the
same retention, since uniformly over subsets
\[
 \sum_i|\rho_i-\be\sigma_i^3|
 \le3\left(R+\be\sum_i\sigma_i^3\right)^{2/3}
                       \left(\sum_iE_i\right)^{1/3}=o(R).
\]
Counting retained variances gives $mR^2/\be^2\to1$.
Replace the core's individually scaled targets by $W_+/\sqrt m$.
Their summed squared scale error tends to zero, and their summed
cubed scale error is $o(R)$; both follow by expanding
$\sigma_i-1/\sqrt m=R(\sigma_i/R-s_0)+R(s_0-1/(R\sqrt m))$.
The outside-to-zero costs are exactly its variance and original price.
This proves the stated IID conclusion after restoring $B$.

In gap units $h=\nus R$, the core's total cubic transport error is
$o(M)$. Remove indices whose error in gap units exceeds a positive number tending
slowly to zero. Their count is $o(M)$ and the preceding uniform moment
continuity estimate preserves variance and price retention. This gives
\eqref{eq:raw-core}; its moment limits follow by uniform $W_3$ continuity.
\end{proof}

\section{A first nonzero Fourier peak and integer lifts}\label{sec:alias}
\begin{lemma}\label{lem:first-alias}
At a gap asymptotic to that in Proposition~\ref{prop:core}, the original
variables in gap units have representations
\[
 Y_i=J_i+Z_i,\qquad J_i=K_i-\E K_i,\quad K_i\in\Z,\quad
 \E Z_i=0,\quad |Z_i|\le1.
\]
The pairs are independent between coordinates. On the retained core
$K_i$ converges uniformly in $W_3$, after integer translations, to
$\Bern(\pe)$. Outside it the summed centered integer variance and
cubic price are $o(M)$. Writing $K=\sum K_i$ and $T=\sum Z_i$,
\begin{equation}\label{eq:initial-lifts}
 Q=\Var T=O(1),\qquad \Cov(K,T)=o(\sqrt M),\qquad
 A=\Var K\sim\pe\qe M.
\end{equation}
\end{lemma}
\begin{proof}
Fix $0<\epsilon<1/2$ in Lemma~\ref{lem:window}. The frequency in gap units
window stays in a compact subinterval of $(0,4\pi)$ containing only
the alias $2\pi$, and its actual characteristic-modulus area is
at least $c/\sqrt M$.
Uniform gap-unit $W_3$ convergence in \eqref{eq:raw-core} gives uniform
convergence of every core characteristic function and its first two
derivatives on that window. Polynomial-factor differences are controlled
by the $L^2$ coupling and the bounded Bernoulli target, so no fourth
moment is required. Around $2\pi$ those factors are nonzero and their
log-modulus second derivatives are at most $-c$. Away from a fixed
neighborhood the retained product is exponentially small.

Let $s_0\to2\pi$ be the exact retained-product modulus maximum.
Its product is bounded by $\exp(-cM(s-s_0)^2)$ near that point.
The area lower bound therefore yields a full-product point
$s_1=s_0+O(M^{-1/2})$ of fixed positive height. The centered omitted
original sum has variance $o(M)$, so its characteristic function changes
by $o(1)$ from $s_1$ to $s_0$:
\[
 |\chi_{\rm rem}(s_1)-\chi_{\rm rem}(s_0)|
 \le |s_1-s_0|\sqrt{\Var Y_{\rm rem}}=o(1).
\]
The full product at the exact core maximum also has fixed positive height.
Rescale the original variables by $s_0/(2\pi)$.

Align each marginal characteristic argument at $2\pi$ and write
$Y_i=\theta_i+K_i+E_i$ with $E_i\in[-1/2,1/2]$, using a fixed
nearest-integer tie rule. Then
\[
 \E\sin(2\pi E_i)=0,\quad \E\cos(2\pi E_i)=r_i>0,\quad
 \sum_i\E E_i^2\le\tfrac18\sum_i(1-r_i)=O(1).
\]
Centering gives $Z_i=E_i-\E E_i$ and the claimed representation.
The exact core logarithmic derivative is zero, hence
$\sum_{I}\E[Y_i\sin(2\pi E_i)]/r_i=0$.
On the core $r_i\to1$ uniformly. Removing these denominators costs
$o(1)$ by Cauchy and the bounded sum of defects. Replacing the sine
by $2\pi E_i$ costs $O(1)+o(\sqrt M)$: couple $Y_i$ to its bounded
binary target and use
\[
 \sum_I\E[|Y_i||E_i|^3]
 \le K\sum_I\E|E_i|^3+
       \left(\sum_I\E|Y_i-W_i|^2\right)^{1/2}
       \left(\sum_I\E|E_i|^6\right)^{1/2}.
\]
The omitted covariance is $o(\sqrt M)$ by its original variance and
bounded phase variance. Thus $\Cov(\sum Y_i,T)=o(\sqrt M)$ and
$\Cov(K,T)=\Cov(\sum Y_i,T)-Q=o(\sqrt M)$.

Core rounding gives uniform $W_3$ convergence of $K_i$ to $\Bern(\pe)$:
outside a fixed small coupling-error event its two nearest integers are
exactly the binary labels, while the remaining cubic rounding error
is bounded by a constant times the cubic target-coupling error.
For the omitted laws, $J_i=Y_i-Z_i$ and $|Z_i|\le1$ imply
\[
 \sum_{\rm rem}\Var K_i\le2\sum_{\rm rem}\Var Y_i+2Q=o(M),\quad
 \sum_{\rm rem}\E|J_i|^3\le4\sum_{\rm rem}\E|Y_i|^3+4Q=o(M).
\]
All prices in the last expression remain separate original-coordinate
prices. These estimates also give $A\sim\pe\qe M$.
\end{proof}
